\Abstract{
Selbstbestimmte Identität (Self-Sovereign Identity) gilt als tragfähiger Ansatz für das dezentrale Management von Gesundheitsdaten. Die dafür nötige Vertrauensverankerung übernimmt in nahezu allen publizierten Architekturen eine Blockchain; hashgraphbasierte Verfahren sind bislang kaum untersucht, und die vorliegenden Evaluationsergebnisse sind wegen unterschiedlicher Messgrößen, Lastprofile und Systemgrenzen untereinander nicht vergleichbar. Ob sich ein reiner Ordnungsdienst ohne eigene Zustandshaltung überhaupt als Verifiable Data Registry eignet, ist damit offen.

Diese Arbeit beantwortet die Frage experimentell. Ausgehend von einer sicherheitsgetrieben entworfenen Referenzarchitektur für den Austausch medizinischer Daten wird deren auf Hyperledger Indy beruhende Ledger-Schicht durch eine auf dem Hedera Consensus Service beruhende Alternative ersetzt, während die darüberliegende Agenten- und Kommunikationsschicht unverändert bleibt. Da Hedera Nachrichten ausschließlich ordnet und ihren Inhalt nicht interpretiert, müssen Zustandshaltung und Berechtigungsprüfung oberhalb des Ledgers rekonstruiert werden. Diese Rekonstruktion wird in sieben dokumentierten Entwurfsentscheidungen entwickelt, als Proof-of-Concept umgesetzt und in fünf Lastszenarien gegen die unveränderte Referenzvariante vermessen: Basismessung, Laststeigerung, Datenwachstum, Knotenausfall und Zugriffskontrolle.

Die funktionale Substitution gelingt vollständig: Sämtliche Kernfunktionen der Ledger-Schicht lassen sich mit den Mechanismen des Ordnungsdienstes nachbilden, und kein Szenario scheitert an einer fehlenden Fähigkeit der Technologie. Die gemessene Leistung fällt dagegen deutlich zugunsten der Indy-Variante aus, in der Latenz ebenso wie im Sättigungsverhalten und in der nutzbaren Fehlertoleranz. Der wesentliche Beitrag dieser Arbeit liegt in der Zuordnung dieser Nachteile: Drei der vier Hauptbefunde lassen sich bis auf die Quellcodeebene der verwendeten Klientbibliotheken zurückverfolgen und sind damit nicht dem Hashgraph-Konsensverfahren zuzuschreiben. Aussagen über eine Technologie und Aussagen über ihren Softwarestand müssen deshalb getrennt gelesen werden. Als Nebenergebnis entstehen vier zuvor nicht gemeldete Einschränkungen dieser Bibliotheken, die ein systematisches Reifegradgefälle zwischen Lese- und Schreibpfad erkennen lassen.
}